\documentclass[10pt,a4paper]{amsart}
\usepackage[margin=19mm]{geometry}
\usepackage{amsmath,amssymb,mathtools}
\usepackage[scr=rsfs,cal=euler]{mathalfa}
\usepackage{xfrac}
\usepackage[unicode,hidelinks,bookmarks=false]{hyperref}
\newcommand{\ie}{i.e.\ }
\newcommand{\cf}{cf.\ }
\newcommand{\resp}{respectively\ }
\newcommand{\Subsection}{\subsection}
\providecommand{\nobreakdash}{\nobreak\hbox{-}}
\setlength{\emergencystretch}{2em}
\multlinegap0pt
\usepackage{amssymb}
\newtheorem{lemma}{Lemma}
\title{On the norms of the multiplication operators between weighted Bergman spaces}
\author{Jianjun Jin}
\date{Source: arXiv:2603.16170v9 (CC BY 4.0)}
\begin{document}
\maketitle

\noindent\emph{Source and licence.} On the norms of the multiplication operators between weighted Bergman spaces. Version arXiv:2603.16170v9 (CC BY 4.0), distributed by arXiv under the Creative Commons Attribution 4.0 International licence, http://creativecommons.org/licenses/by/4.0/, as recorded in the arXiv licence metadata for this exact version and shown on the abstract page https://arxiv.org/abs/2603.16170v9. Full licence notice: \texttt{licenses/arxiv-2603.16170-CC-BY-4.0.txt}.\par\smallskip

\noindent\textit{Editorial scope.} Counted unit: the article's lemma l-2 (source0.tex lines 196--207). The article's complete original proof of it is delivered with it (source0.tex lines 210--280, one \texttt{proof} environment of 71 lines, which the article places away from the statement). No mathematics is written, repaired or re-derived: the statement and the proof are the article's own lines, in the article's own notation, and no retained line is substituted or re-flowed.  No further source-local context is needed: every cross-reference of every retained line resolves inside the two delivered spans.  Nothing was omitted from any retained span, so the package carries no omission record.  No retained line contains a citation, so the wrapper carries no bibliography and no bibliography entry.  No retained line contains a figure environment, an image-inclusion command or drawing code, so no image is delivered and the package declares no asset dependency.  Editorial formatting only, in addition: an AMS \texttt{amsart} preamble that restates the article's own declarations and macros used by the retained lines, namely the environments \texttt{lemma} on separate counters, together with the standard packages those lines need.  The retained environments print in this document as Lemma 1 in order of appearance, whereas the article prints the same items with the numbering of its own full document.  The complete native source of the article as distributed in the arXiv e-print for this exact version is bundled unchanged as \texttt{sources/196520/source0.tex}.
\begin{lemma}\label{l-2}
Let $\alpha>-1, r\in (0,1), \lambda>0$. If $2\lambda>\alpha+2$, then  
\begin{equation}\label{a-e-1}
 \|{(1-rz^2)^{-\lambda}}\|_{\alpha}^2=\frac{\Gamma(\alpha+2)\Gamma(2\lambda-\alpha-2)}{2^{\alpha+1}[\Gamma(\lambda)]^2}\cdot\frac{1+{o}(1)}{(1-r^2)^{2\lambda-\alpha-2}},\, {\textup{as}}\,\, r\rightarrow 1^{-}. 
 \end{equation}
 and \begin{equation}\label{a-e-2}
   \|{(1-rz)^{-\lambda}}\|_{\alpha}^2=\frac{\Gamma(\alpha+2)\Gamma(2\lambda-\alpha-2)}{[\Gamma(\lambda)]^2}\cdot\frac{1+{o}(1)}{(1-r^2)^{2\lambda-\alpha-2}},\, {\textup{as}}\,\, r\rightarrow 1^{-}. 
 \end{equation} 
 If $0<2\lambda<\alpha+2$, then there is a constant $M>0$ such that for all $r\in (0,1)$,
 \begin{equation}\label{a-e-3}
   \|{(1-rz)^{-\lambda}}\|_{\alpha}^2\leq M.
 \end{equation}\end{lemma} 

\begin{proof}[Proof of Lemma \ref{l-2}]
First, for $\gamma>0$, we have 
\begin{equation}\label{l-eq-1}
  \frac{1}{(1-z)^{\gamma}}=\sum_{n=0}^{\infty}\frac{\Gamma(n+\gamma)}{n!\Gamma(\gamma)}z^n, z\in \Delta.\end{equation}
Also, for any $\phi=\sum_{n=0}^{\infty}a_nz^n\in \mathbf{A}_{\alpha}^2(\Delta)$, we have
\begin{equation}\label{l-eq-2}\|\phi\|_{\alpha}^2=\sum_{n=0}^{\infty}\frac{n!\Gamma(\alpha+2)}{\Gamma(n+\alpha+2)}|a_n|^{2}.
\end{equation}
For $r\in (0,1)$ and $\lambda>0$, from (\ref{l-eq-1}) and (\ref{l-eq-2}), we obtain that
  \begin{equation}\label{l-eq-3}
 \|{(1-rz^2)^{-\lambda}}\|_{\alpha}^2=\sum_{n=0}^{\infty}\frac{(2n)!\Gamma(\alpha+2)}{\Gamma(2n+\alpha+2)}\Big|\frac{\Gamma(n+\lambda)}{n!\Gamma(\lambda)}\Big|^2 r^{2n}, 
 \end{equation}
and 
 \begin{equation}\label{l-eq-4}
 \|{(1-rz)^{-\lambda}}\|_{\alpha}^2=\sum_{n=0}^{\infty}\frac{n!\Gamma(\alpha+2)}{\Gamma(n+\alpha+2)}\Big|\frac{\Gamma(n+\lambda)}{n!\Gamma(\lambda)}\Big|^2 r^{2n}.
 \end{equation}
By Stirling's formula, we have
 \begin{equation}
 \frac{\Gamma(n+\lambda)}{n!}=(n+1)^{\lambda-1}[1+\widetilde{o}(1)],\,\, {\textup{as}} \,\, n\rightarrow \infty.\nonumber 
 \end{equation}
 Here and later, we use $\widetilde{o}(1)$ to denote a sequence that tends to $0$, i.e., $o(1)\rightarrow 0, \, {\textup{as}}\,\, n \rightarrow \infty$, which may be different in different places. Then it follows from (\ref{l-eq-3}) that 
 \begin{eqnarray}\label{l-eq-5}
 \|{(1-rz^2)^{-\lambda}}\|_{\alpha}^2&=&\sum_{n=0}^{\infty}\frac{\Gamma(\alpha+2)}{[2(n+1)]^{\alpha+1}}[1+\widetilde{o}(1)]\cdot\frac{(n+1)^{2\lambda-2}}{[\Gamma(\lambda)]^2}[1+\widetilde{o}(1)]r^{2n}\nonumber \\
 &=&\frac{1}{2^{\alpha+1}}\frac{\Gamma(\alpha+2)}{[\Gamma(\lambda)]^2}\sum_{n=0}^{\infty}{(n+1)^{2\lambda-\alpha-3}}[1+\widetilde{o}(1)]r^{2n}.
 \end{eqnarray}
Similarly, from (\ref{l-eq-4}), we have 
 \begin{eqnarray}\label{l-eq-6}
 \|{(1-rz)^{-\lambda}}\|_{\alpha}^2=\frac{\Gamma(\alpha+2)}{[\Gamma(\lambda)]^2}\sum_{n=0}^{\infty}{(n+1)^{2\lambda-\alpha-3}}[1+\widetilde{o}(1)]r^{2n}.
 \end{eqnarray}
Note that there is a constant $C>0$ such that for $a<-1$,
$$\sum_{n=1}^{\infty}n^a r^{2n}\leq C,\,\,{\text{for all}}\,\, r\in (0,1).$$
Consequently, when $0<2\lambda<\alpha+2$ so that $2\lambda-\alpha-3<-1$, we see that (\ref{a-e-3}) holds. 

On the other hand, when $2\lambda>\alpha+2$, we obtain from Stirling's formula and \ref{l-eq-1} that 
  \begin{eqnarray}\label{l-eq-7}
 \lefteqn{\sum_{n=0}^{\infty}{(n+1)^{2\lambda-\alpha-3}}[1+\widetilde{o}(1)]r^{2n}=\sum_{n=0}^{\infty}\frac{\Gamma(n+2\lambda-\alpha-2)}{n!\Gamma(2\lambda-\alpha-2)}[1+\widetilde{o}(1)]r^{2n}}\nonumber \\
 &&=\frac{1}{{(1-r^2)^{2\lambda-\alpha-2}}}+\sum_{n=0}^{\infty}\frac{\Gamma(n+2\lambda-\alpha-2)}{n!\Gamma(2\lambda-\alpha-2)}\cdot [\widetilde{o}(1)]\cdot r^{2n}. 
 \end{eqnarray}
Let $\widetilde{o}(1):=b_n$. Note that for any $\varepsilon\in (0,1/2)$ there exist two constants $C_0>0, N_0\in \mathbb{N}$ such that $|b_n|\leq C_0$ for all $n\in \mathbb{N}_0$ and $|b_n|\leq \varepsilon/2$ for all $n\geq N_0$ in the last term of (\ref{l-eq-7}). 
It follows from again  (\ref{l-eq-1}) that
\begin{eqnarray}\label{l-eq-8}
\lefteqn{\sum_{n=0}^{\infty}\frac{\Gamma(n+2\lambda-\alpha-2)}{n!\Gamma(2\lambda-\alpha-2)}\cdot [\widetilde{o}(1)]\cdot r^{2n}} \nonumber \\
  && \leq C_0\sum_{n=0}^{N_0} \frac{\Gamma(n+2\lambda-\alpha-2)}{n!\Gamma(2\lambda-\alpha-2)}r^{2n}+\varepsilon/2\sum_{n={{n_0}+1}}^{\infty}\frac{\Gamma(n+2\lambda-\alpha-2)}{n!\Gamma(2\lambda-\alpha-2)}r^{2n} \nonumber \\
  && \leq C_0\sum_{n=0}^{N_0} \frac{\Gamma(n+2\lambda-\alpha-2)}{n!\Gamma(2\lambda-\alpha-2)}+\frac{\varepsilon/2}{{(1-r^2)^{2\lambda-\alpha-2}}}\nonumber \\
  &&:=C_1+\frac{\varepsilon/2}{{(1-r^2)^{2\lambda-\alpha-2}}}.
  \end{eqnarray}
Meanwhile, \begin{eqnarray}\label{l-eq-9}
\lefteqn{\sum_{n=0}^{\infty}\frac{\Gamma(n+2\lambda-\alpha-2)}{n!\Gamma(2\lambda-\alpha-2)}\cdot [\widetilde{o}(1)]\cdot r^{2n}} \nonumber \\
  && \geq \sum_{n=0}^{N_0} \frac{\Gamma(n+2\lambda-\alpha-2)}{n!\Gamma(2\lambda-\alpha-2)}b_n r^{2n}-\varepsilon/2\sum_{n={{n_0}+1}}^{\infty}\frac{\Gamma(n+2\lambda-\alpha-2)}{n!\Gamma(2\lambda-\alpha-2)}r^{2n} \nonumber \\
  && \geq -\sum_{n=0}^{N_0} \frac{\Gamma(n+2\lambda-\alpha-2)}{n!\Gamma(2\lambda-\alpha-2)}|b_n|-\frac{\varepsilon/2}{{(1-r^2)^{2\lambda-\alpha-2}}}\nonumber \\
  &&:=-C_2-\frac{\varepsilon/2}{{(1-r^2)^{2\lambda-\alpha-2}}}.
  \end{eqnarray}
Combining (\ref{l-eq-7}), (\ref{l-eq-8}) and (\ref{l-eq-9}), we obtain that
    \begin{eqnarray}
[{(1-r^2)^{2\lambda-\alpha-2}}]\sum_{n=0}^{\infty}{(n+1)^{2\lambda-\alpha-3}}[1+\widetilde{o}(1)]r^{2n}\leq 1+\varepsilon/2+C_1[{(1-r^2)^{2\lambda-\alpha-2}}], \nonumber
 \end{eqnarray}
 and 
     \begin{eqnarray}
[{(1-r^2)^{2\lambda-\alpha-2}}]\sum_{n=0}^{\infty}{(n+1)^{2\lambda-\alpha-3}}[1+\widetilde{o}(1)]r^{2n}\geq 1-\varepsilon/2-C_2[{(1-r^2)^{2\lambda-\alpha-2}}]. \nonumber
 \end{eqnarray}
Since $C_1$ and $C_2$ are independent of $r$, then we can find a constant $\delta\in (0,1)$ such that
$$C_1{(1-r^2)^{2\lambda-\alpha-2}}<\varepsilon/2,\,\,\text{and}\,\, C_2{(1-r^2)^{2\lambda-\alpha-2}}<\varepsilon/2,$$
for all $r\in (1-\delta, 1)$. Consequently, for any $\varepsilon\in (0,1/2)$, we see that 
    \begin{eqnarray}
1-\varepsilon\leq [{(1-r^2)^{2\lambda-\alpha-2}}]\sum_{n=0}^{\infty}{(n+1)^{2\lambda-\alpha-3}}[1+\widetilde{o}(1)]r^{2n}\leq 1+\varepsilon, \nonumber
 \end{eqnarray}
for all $r\in (1-\delta, 1)$. This is equivalent to  
 \begin{equation}\label{l-eq-10}
\sum_{n=0}^{\infty}{(n+1)^{2\lambda-\alpha-3}}[1+\widetilde{o}(1)]r^{2n}=\frac{1+{o}(1)}{(1-r^2)^{2\lambda-\alpha-2}},\, {\textup{as}}\,\, r\rightarrow 1^{-}. 
 \end{equation}
 Then, in view of (\ref{l-eq-10}), (\ref{a-e-1}) and (\ref{a-e-2}) respectively follow from (\ref{l-eq-5}) and (\ref{l-eq-6}).  The lemma is proved.
 \end{proof} 

\end{document}
